\scanpage{16}
\mspage{7}
\uline{Proposition} If $(G=(S,\widetilde A,\sigma),I,\uline{g}:S\unsure{\to|-}\mathbb N)$ is an MD-graph, its type $(g,\nu)$ is anabelian, i.e. $2g+\nu\geq3$.

If we had $g=1$, $\nu=0$, then the relation
\[
g=1=\sum g_\alpha+h_{\unsure{1|0}}
\]
shows that either all the $g_\alpha$ are zero and $h_1=1\deleted{\unsure{+1|and}}$, or all the $g_\alpha$ except one \inserted{$g_{\alpha_0}$} are zero, and $g_{\alpha_0}=1$, $h_1=0$ \inserted{(i.e. $G$ is a tree)}. \inserted{In the first case,} The relation $s_0-s_1=h_0-h_1=1-1=0$ shows that there are as many vertices as edges, in the second, $s_0-s_1=1-0=1$ there is one more vertex than edges.

But the relation $\nu=0$ implies that the vertices of genus $0$ are of degree $\geq3$. \inserted{$G$ is not reduced to a vertex, since it is stable of type $(1,0)$\unsure{.}} But a \deleted{compact topological graph \unsure{or}} \inserted{tree which is not reduced to a single vertex} has at least two ends, of which at least one will be of genus $0$, which eliminates case 2°. Thus all the vertices are of genus zero \inserted{(thus no vertex is an end)} and $h_1=\unsure{1|0}$. But \deleted{\unsure{since}} a $0$-connected topological graph with $h_1=1$ which has no \inserted{point of} topological boundary is homeomorphic to the circle, so if it is realized by a combinatorial graph, all the vertices of the latter are of degree $2$, which contradicts their being of degree three.
\mspage{15}
